% Môn: Vật lý
% Lớp: 12
% Chương: Chương I. Vật lí nhiệt
% Bài: L12C1 Bài 4. Nhiệt dung riêng · Thực nghiệm xác định nhiệt hóa hơi riêng
% ===== Câu 81 =====
% ID: L12C1B4-142-DS
% Mức: NB
\dangbt{Thực nghiệm xác định nhiệt hóa hơi riêng}
\begin{ex}
% ID: L12C1B4-142-DS
% Mức: NB
Về xử lí số liệu.
\choiceTF
{\True Có thể tính L_i cho từng lần đo rồi lấy trung bình.}
{\True Cần ghi đơn vị và số chữ số có nghĩa hợp lí.}
{Nên loại tùy ý mọi số liệu không vừa ý.}
{\True Có thể dùng độ dốc đồ thị m theo t để xác định tốc độ hóa hơi.}
\loigiai{
\begin{itemchoice}
\item (Đúng) Có thể tính L_i cho từng lần đo rồi lấy trung bình.
\item (Đúng) Cần ghi đơn vị và số chữ số có nghĩa hợp lí.
\item (Sai) Nên loại tùy ý mọi số liệu không vừa ý.
\item (Đúng) Có thể dùng độ dốc đồ thị m theo t để xác định tốc độ hóa hơi.
\end{itemchoice}
}
\end{ex}

% ===== Câu 85 =====
% ID: L12C1B4-143-DS
% Mức: TH
\begin{ex}
% ID: L12C1B4-143-DS
% Mức: TH
Các nguồn sai số trong thí nghiệm.
\choiceTF
{\True Nhiệt truyền ra môi trường.}
{\True Hơi nước ngưng tụ rồi rơi lại vào ấm.}
{\True Gió làm thay đổi tốc độ mất nhiệt và hóa hơi.}
{\True Cân có độ chia hữu hạn.}
\loigiai{
\begin{itemchoice}
\item (Đúng) Nhiệt truyền ra môi trường.
\item (Đúng) Hơi nước ngưng tụ rồi rơi lại vào ấm.
\item (Đúng) Gió làm thay đổi tốc độ mất nhiệt và hóa hơi. (Vì): ả bốn yếu tố đều có thể ảnh hưởng kết quả
\item (Đúng) Cân có độ chia hữu hạn.
\end{itemchoice}
}
\end{ex}

% ===== Câu 89 =====
% ID: L12C1B4-144-DS
% Mức: VD
\begin{ex}
% ID: L12C1B4-144-DS
% Mức: VD
Thiết kế thí nghiệm xác định $L$ bằng ấm điện.
\choiceTF
{\True Cần xác định công suất điện hoặc năng lượng cung cấp.}
{\True Cần đo khối lượng nước hóa hơi.}
{\True Nên lấy số liệu khi nước sôi ổn định.}
{Có thể bỏ qua mọi hao phí mà không ảnh hưởng kết quả.}
\loigiai{
\begin{itemchoice}
\item (Đúng) Cần xác định công suất điện hoặc năng lượng cung cấp.
\item (Đúng) Cần đo khối lượng nước hóa hơi. (Vì): a đại lượng đầu cần thiết; hao phí là nguồn sai số hệ thống
\item (Đúng) Nên lấy số liệu khi nước sôi ổn định.
\item (Sai) Có thể bỏ qua mọi hao phí mà không ảnh hưởng kết quả.
\end{itemchoice}
}
\end{ex}

% ===== Câu 93 =====
% ID: L12C1B4-145-DS
% Mức: VD
\begin{ex}
% ID: L12C1B4-145-DS
% Mức: VD
Đồ thị m(t) trong giai đoạn sôi ổn định.
\choiceTF
{\True Khối lượng còn lại giảm theo thời gian.}
{\True Nếu tốc độ hóa hơi không đổi, đồ thị gần đường thẳng.}
{\True Độ lớn hệ số góc là tốc độ giảm khối lượng.}
{Đồ thị luôn tăng vì nhiệt độ tăng.}
\loigiai{
\begin{itemchoice}
\item (Đúng) Khối lượng còn lại giảm theo thời gian.
\item (Đúng) Nếu tốc độ hóa hơi không đổi, đồ thị gần đường thẳng.
\item (Đúng) Độ lớn hệ số góc là tốc độ giảm khối lượng.
\item (Sai) Đồ thị luôn tăng vì nhiệt độ tăng.
\end{itemchoice}
}
\end{ex}

% ===== Câu 97 =====
% ID: L12C1B4-146-DS
% Mức: VDC
\begin{ex}
% ID: L12C1B4-146-DS
% Mức: VDC
Nguồn hữu ích 800 $W$ làm bay hơi $0,040\,\text{kg}$ trong $100\,\text{s}$.
\choiceTF
{\True Năng lượng hữu ích là 80 kJ.}
{\True $L$ tính được là 2,0.10^ $6\,\text{J}$ /kg.}
{\True Nếu khối lượng đọc thiếu thì $L$ tính được có xu hướng lớn.}
{Đơn vị của $L$ là W.}
\loigiai{
\begin{itemchoice}
\item (Đúng) Năng lượng hữu ích là 80 kJ.
\item (Đúng) $L$ tính được là 2,0.10^ $6\,\text{J}$ /kg.
\item (Đúng) Nếu khối lượng đọc thiếu thì $L$ tính được có xu hướng lớn.
\item (Sai) Đơn vị của $L$ là W.
\end{itemchoice}
}
\end{ex}

% ===== Câu 98 =====
% ID: L12C1B4-147-TL
% Mức: NB
\begin{bt}
% ID: L12C1B4-147-TL
% Mức: NB
Nêu ít nhất bốn biện pháp giảm sai số khi xác định L.
\loigiai{
Đợi sôi ổn định mới đo; che chắn gió nhưng bảo đảm hơi thoát an toàn; dùng cân có độ phân giải phù hợp; chọn khoảng thời gian đủ dài để Δm đủ lớn; đo công suất thực; hạn chế hơi ngưng tụ rơi lại; lặp nhiều lần và lấy trung bình.
}
\end{bt}

% ===== Câu 99 =====
% ID: L12C1B4-148-TLN
% Mức: NB
\begin{ex}
% ID: L12C1B4-148-TLN
% Mức: NB
Ấm 1200 $W$, hiệu suất hữu ích 75%, trong $150\,\text{s}$ làm bay hơi $60\,\text{g}$. Tính $L$ theo đơn vị 10^ $6\,\text{J}$ /kg, chỉ ghi số.
\shortans{2.25}
\loigiai{
$L=0,75.1200.150/0,060=2,25.10^6 J/kg.$
}
\end{ex}

% ===== Câu 102 =====
% ID: L12C1B4-149-TL
% Mức: TH
\begin{bt}
% ID: L12C1B4-149-TL
% Mức: TH
Trình bày phương án thí nghiệm xác định nhiệt hóa hơi riêng của nước bằng ấm điện, cân và đồng hồ.
\loigiai{
Đặt ấm có nước trên cân, đun đến sôi ổn định. Ghi khối lượng đầu m1, bắt đầu bấm giờ; sau thời gian t ghi m2. Khối lượng hóa hơi Δm=m1-m2. Xác định công suất hữu ích P_h hoặc công suất $P$ và hiệu suất H. Tính L=P_ht/Δm=HPt/Δm. Lặp lại và lấy trung bình.
}
\end{bt}

% ===== Câu 103 =====
% ID: L12C1B4-150-TLN
% Mức: NB
\begin{ex}
% ID: L12C1B4-150-TLN
% Mức: NB
Đo được $L$ lần lượt 2,20; 2,30; 2,25 (đơn vị 10^ $6\,\text{J}$ /kg). Giá trị trung bình theo cùng đơn vị là bao nhiêu? Chỉ ghi số.
\shortans{2.25}
\loigiai{
$L_tb=(2,20+2,30+2,25)/3=2,25.$
}
\end{ex}

% ===== Câu 106 =====
% ID: L12C1B4-151-TL
% Mức: TH
\begin{bt}
% ID: L12C1B4-151-TL
% Mức: TH
Từ đồ thị m(t)=0,800-0,00040t (kg), công suất hữu ích 920 $W$, xác định L.
\loigiai{
Độ lớn tốc độ hóa hơi k= $0,00040\,\text{kg}$ /s. Vì P_h=Lk nên L=P_h/k=920/0,00040=2,30.10^ $6\,\text{J}$ /kg.
}
\end{bt}

% ===== Câu 107 =====
% ID: L12C1B4-152-TLN
% Mức: TH
\begin{ex}
% ID: L12C1B4-152-TLN
% Mức: TH
Trong giai đoạn sôi, khối lượng giảm từ $850\,\text{g}$ xuống $790\,\text{g}$. Khối lượng hóa hơi theo gam là bao nhiêu? Chỉ ghi số.
\shortans{60}
\loigiai{
Δm=850-790= $60\,\text{g}$.
}
\end{ex}

% ===== Câu 110 =====
% ID: L12C1B4-153-TL
% Mức: VD
\begin{bt}
% ID: L12C1B4-153-TL
% Mức: VD
Một ấm 1000 $W$, hiệu suất 80%. Khi nước sôi ổn định, sau $240\,\text{s}$ khối lượng giảm $84\,\text{g}$. Tính L.
\loigiai{
Năng lượng hữu ích Q=0,80.1000.240= $192000\,\text{J}$. Δm= $0,084\,\text{kg}$. L=192000/0,084≈2,29.10^ $6\,\text{J}$ /kg.
}
\end{bt}

% ===== Câu 111 =====
% ID: L12C1B4-154-TLN
% Mức: VD
\begin{ex}
% ID: L12C1B4-154-TLN
% Mức: VD
Đồ thị cho khối lượng giảm đều $24\,\text{g}$ trong $80\,\text{s}$. Tốc độ giảm khối lượng theo g/s là bao nhiêu? Chỉ ghi số.
\shortans{0.3}
\loigiai{
Tốc độ |Δm/Δt|=24/80= $0,30\,\text{g}$ /s.
}
\end{ex}

% ===== Câu 114 =====
% ID: L12C1B4-155-TL
% Mức: VD
\begin{bt}
% ID: L12C1B4-155-TL
% Mức: VD
Giải thích vì sao phải phân biệt công suất điện định mức và công suất hữu ích khi xử lí thí nghiệm.
\loigiai{
Công suất điện định mức mô tả năng lượng thiết bị tiêu thụ mỗi giây; chỉ một phần truyền cho nước để hóa hơi. Phần còn lại làm nóng ấm và thất thoát. Dùng công suất định mức như công suất hữu ích sẽ làm sai $L$; cần đo/ước lượng hiệu suất hoặc công suất hao phí.
}
\end{bt}

% ===== Câu 115 =====
% ID: L12C1B4-156-TLN
% Mức: VDC
\begin{ex}
% ID: L12C1B4-156-TLN
% Mức: VDC
Công suất hữu ích 900 $W$ làm bay hơi $45\,\text{g}$ trong $100\,\text{s}$. Tính $L$ theo đơn vị 10^ $6\,\text{J}$ /kg, chỉ ghi số.
\shortans{2}
\loigiai{
$L=900.100/0,045=2,0.10^6 J/kg.$
}
\end{ex}

% ===== Câu 118 =====
% ID: L12C1B4-157-TL
% Mức: VDC
\begin{bt}
% ID: L12C1B4-157-TL
% Mức: VDC
Phân tích ảnh hưởng của thất thoát nhiệt nếu dùng L=Pt/m mà không hiệu chỉnh.
\loigiai{
Pt gồm năng lượng hóa hơi Lm và năng lượng thất thoát Q_hp. Do đó Pt/m=L+Q_hp/m>L. Kết quả bị lệch lớn theo chiều tăng. Cần xác định hiệu suất/công suất hữu ích hoặc hiệu chỉnh công suất hao phí.
}
\end{bt}

% ===== Câu 119 =====
% ID: L12C1B4-158-TLN
% Mức: VDC
\begin{ex}
% ID: L12C1B4-158-TLN
% Mức: VDC
Công suất hữu ích 690 $W, L=2,3.10^6\,\text{J}$ /kg. Trong $200\,\text{s}$, khối lượng nước hóa hơi theo gam là bao nhiêu? Chỉ ghi số.
\shortans{60}
\loigiai{
$m=P_ht/L=690.200/(2,3.10^6)=0,060 kg=60 g.$
}
\end{ex}

% ===== Câu 144 =====
% ID: L12C1B4-159-TN
% Mức: NB
\begin{ex}
% ID: L12C1B4-159-TN
% Mức: NB
Trong thí nghiệm dùng ấm điện xác định $L$, đại lượng cần đo trực tiếp thường gồm
\choice
{\True khối lượng nước bay hơi và thời gian đun}
{chỉ nhiệt độ phòng}
{chỉ thể tích ấm}
{chỉ áp suất khí quyển}
\loigiai{
Với công suất biết trước, đo thời gian và độ giảm khối lượng để suy ra L.
}
\end{ex}

% ===== Câu 145 =====
% ID: L12C1B4-160-TN
% Mức: NB
\begin{ex}
% ID: L12C1B4-160-TN
% Mức: NB
Nếu bỏ qua thất thoát nhiệt trong khi thực tế có hao phí, giá trị $L$ tính bằng Pt/m thường
\choice
{nhỏ hơn thực}
{bằng 0}
{\True lớn hơn thực}
{luôn chính xác}
\loigiai{
Pt gồm cả năng lượng thất thoát; gán toàn bộ cho hóa hơi làm $L$ bị tính lớn.
}
\end{ex}

% ===== Câu 146 =====
% ID: L12C1B4-161-TN
% Mức: NB
\begin{ex}
% ID: L12C1B4-161-TN
% Mức: NB
Sai số của khối lượng hóa hơi tương đối lớn khi
\choice
{\True độ giảm khối lượng quá nhỏ so với độ chia cân}
{độ giảm khối lượng đủ lớn}
{cân được đặt ổn định}
{đo lặp lại}
\loigiai{
Đại lượng đo gần độ phân giải của cân tạo sai số tương đối lớn.
}
\end{ex}

% ===== Câu 147 =====
% ID: L12C1B4-162-TN
% Mức: TH
\begin{ex}
% ID: L12C1B4-162-TN
% Mức: TH
Nếu công suất hữu ích P_h đã biết, công thức lí tưởng xác định $L$ là
\choice
{$L=m/(P_ht)$}
{\True $L=P_ht/m$}
{$L=P_hm/t$}
{$L=t/(P_hm)$}
\loigiai{
Năng lượng hữu ích P_ht bằng Lm, nên L=P_ht/m.
}
\end{ex}

% ===== Câu 148 =====
% ID: L12C1B4-163-TN
% Mức: TH
\begin{ex}
% ID: L12C1B4-163-TN
% Mức: TH
Biện pháp giúp giảm sai số ngẫu nhiên là
\choice
{chỉ đo một lần thật nhanh}
{\True đo nhiều lần và lấy trung bình}
{bỏ ghi đơn vị}
{dùng khối lượng rất nhỏ}
\loigiai{
Lặp phép đo và lấy trung bình giúp giảm sai số ngẫu nhiên.
}
\end{ex}

% ===== Câu 149 =====
% ID: L12C1B4-164-TN
% Mức: TH
\begin{ex}
% ID: L12C1B4-164-TN
% Mức: TH
Đồ thị khối lượng nước m còn lại theo thời gian trong giai đoạn sôi ổn định gần là
\choice
{đường tăng dần}
{\True đường giảm gần tuyến tính}
{đường nằm ngang tuyệt đối}
{parabol luôn đi qua gốc}
\loigiai{
Với công suất hữu ích gần không đổi, tốc độ mất khối lượng gần không đổi.
}
\end{ex}

% ===== Câu 150 =====
% ID: L12C1B4-165-TN
% Mức: VD
\begin{ex}
% ID: L12C1B4-165-TN
% Mức: VD
Cân điện tử trong thí nghiệm dùng để đo
\choice
{công suất}
{\True độ giảm khối lượng do nước hóa hơi}
{nhiệt hóa hơi riêng trực tiếp}
{hiệu điện thế}
\loigiai{
Khối lượng hóa hơi là độ giảm số chỉ của cân trong giai đoạn sôi ổn định.
}
\end{ex}

% ===== Câu 151 =====
% ID: L12C1B4-166-TN
% Mức: VD
\begin{ex}
% ID: L12C1B4-166-TN
% Mức: VD
Độ giảm khối lượng $0,050\,\text{kg}$ trong $100\,\text{s}$ với công suất hữu ích 1000 $W$ cho $L$ bằng
\choice
{$2,0.10^4 J/kg$}
{$2,0.10^5 J/kg$}
{\True $2,0.10^6 J/kg$}
{$5,0.10^6 J/kg$}
\loigiai{
$L=P_ht/m=1000.100/0,050=2,0.10^6 J/kg.$
}
\end{ex}

% ===== Câu 152 =====
% ID: L12C1B4-167-TN
% Mức: VDC
\begin{ex}
% ID: L12C1B4-167-TN
% Mức: VDC
Để giảm ảnh hưởng của giai đoạn tăng nhiệt độ, nên bắt đầu tính thời gian khi
\choice
{vừa bật nguồn}
{\True nước đã sôi ổn định}
{nước còn lạnh}
{đã tắt nguồn}
\loigiai{
Khi sôi ổn định, năng lượng chủ yếu dành cho hóa hơi.
}
\end{ex}

% ===== Câu 153 =====
% ID: L12C1B4-168-TN
% Mức: VDC
\begin{ex}
% ID: L12C1B4-168-TN
% Mức: VDC
Trong phương pháp so sánh hai khoảng thời gian ở cùng chế độ sôi, mục đích chính là
\choice
{\True tách ảnh hưởng công suất hao phí gần ổn định}
{làm nước đông đặc}
{đổi đơn vị tự động}
{làm $L$ bằng 0}
\loigiai{
Theo dõi khối lượng theo thời gian khi sôi ổn định giúp ước lượng tốc độ hóa hơi và hiệu chỉnh nền hao phí.
}
\end{ex}
